\documentclass[11pt]{article}


\setlength{\parindent}{0pt}
\usepackage{xltxtra}
\usepackage{hyperref}
\hypersetup{hidelinks}
\usepackage{url}
\urlstyle{tt}
\usepackage{xcolor}
\definecolor{CVBlue}{RGB}{23,110,191}
\usepackage{calc}
\usepackage{graphicx}
\usepackage{tikz}
\usetikzlibrary{calc}
\usepackage{fontspec}
\usepackage{xeCJK}
\usepackage{enumitem}
\CJKsetecglue{} %% 取消中文与数字之间的间隙


%% 主文档字体设置
\setmainfont[
	Path = fonts/Main/,
	Extension = .otf,
	BoldFont = texgyretermes-bold.otf, % 加粗字体
]{texgyretermes-regular.otf} % 正文字体

% 中文字体设置
\setCJKmainfont[
	Path = fonts/hansans/,
	Extension = .ttf,
	BoldFont = NotoSansSC-Bold.ttf, % 加粗字体
]{NotoSansSC-Regular.otf} % 正文字体


\usepackage{relsize}
\usepackage{xspace}

% 使用 fontawesome（部分图标）
\usepackage{fontawesome} 

% A4纸，上下左右边距
\usepackage[
	a4paper,
	left=1.2cm,
	right=1.2cm,
	top=1.5cm,
	bottom=1cm,
	nohead
]{geometry}

\renewcommand{\baselinestretch}{1.5} % 行间距设为1.5

\usepackage{titlesec}
\usepackage{enumitem}
\setlist{noitemsep} % 取消列表项间的额外间距
%\setlist{nosep} % 取消所有垂直间距
\setlist[itemize]{topsep=0.25em, leftmargin=*}
\setlist[enumerate]{topsep=0.25em, leftmargin=*}



\titleformat{\section}         % 定制 \section 命令 
{\large\bfseries\raggedright} % 将 section 标题设置为大号、粗体且左对齐
{}{0em}                      % 可用于为所有 section 添加前缀（如“章节...”）
{}                           % 可用于在标题前插入代码
[{\color{CVBlue}\titlerule}]  % 在标题后插入一条横线
\titlespacing*{\section}{0cm}{*1.6}{*1.2}



\begin{document}
\pagenumbering{gobble}

%%%% 利用tikz来定位照片
\begin{tikzpicture}[remember picture, overlay] 
	\node[anchor = north east] at ($(current page.north east)+(-2cm,-1.2cm)$) {\includegraphics[height=3cm]{avatar.jpg}};
  \end{tikzpicture}%
  %%%% 利用tikz来定位学校Logo，这里只在第一页显示
  \begin{tikzpicture}[remember picture, overlay] 
	\node[anchor = north west] at ($(current page.north west)+(0.5cm,+1.0cm)$) {\includegraphics[height=6cm]{zju.png}};
  \end{tikzpicture}%
\centerline{\LARGE\bfseries{胡豆}} 

\centerline{\normalsize{\faPhone\ 111-1111-1111 \quad \faEnvelopeO\ \href{mailto:3230100000@zju.edu.cn}{3230100000@zju.edu.cn}}} 

\centerline{\normalsize{\faGithubSquare\ \href{https://github.com/maksymilan}{https://github.com/maksymilan} \quad \faRssSquare\ \href{https://maksymilan.github.io/}{https://maksymilan.github.io}}} 
    
\section{\makebox[\widthof{\faGraduationCap}][c]{\color{CVBlue}\faGraduationCap}\ 教育背景}    
\textbf{浙江大学} \hfill 2023.9 -- 至今\\ 
CC98烂坑挖掘工程与技术\quad 大二 
\begin{itemize}[nosep]
    \item 相关课程：《烂坑挖掘及基础》、《高级挖坑技巧》、《烂坑数理统计》
\end{itemize}

\section{\makebox[\widthof{\faUsers}][c]{\color{CVBlue}\faUsers}\ 项目经历}


\textbf{98烂坑风控与可视化作战平台} \hfill 2024.08 -- 至今 \\
项目描述：为维护cc98社区的讨论纯度、净化“心灵之约”版块风气，本项目旨在对“烂坑”进行科学、严谨的量化分析，并引入\textbf{大语言模型(LLM)}对烂坑帖进行\textbf{自动化情感极性与“离谱指数”标注}。
\begin{itemize}[nosep]
    \item \textbf{数据与后端}：运用 \textbf{Go} 语言及 \textbf{GORM} 框架，实现对全站“烂坑”帖的7x24小时全天候\textbf{抓取}与\textbf{入库}(MySQL)。
    \item \textbf{可视化前端}：采用原生 \textbf{HTML/CSS/JS} 构建了“烂坑”实时监控\textbf{作战大屏}。实现了坑主成分分析、挖坑热力图、年度烂坑王榜单等\textbf{可视化图表}，并通过 \textbf{AJAX} 异步刷新，确保数据洞察的实时性。
    \item \textbf{核心产出}：全面掌握了从\textbf{数据采集、清洗、LLM标注到可视化分析}的“反烂坑”全链路技术。
\end{itemize}

\vspace{0.3em} % 

\textbf{AIGC驱动的超进化自动挖坑机 (Agent)} \hfill 2025.02 -- 至今 \\
项目描述：为探索心灵版规的边界，以及测试版主锁沉反应速度的极限，本项目基于\textbf{多智能体协作(MCP)与LLM}技术，开发了一款能\textbf{自主“挖坑”与“互坑”}的AI Agent。
\begin{itemize}[nosep]
    \item \textbf{并发与调度}：后端采用 \textbf{Go} 语言，充分利用 \textbf{Goroutine} 和 \textbf{Channel} 的高并发模型，模拟\textbf{用户并发在线}，执行定时挖坑，确保挖坑行动的\textbf{隐蔽性}与\textbf{高效性}。
    \item \textbf{实时监控与响应}：运用 \textbf{WebSocket} 协议实现了对目标帖子的\textbf{实时监听与交互}。确保在烂坑被识破的第一时间，Agent能光速完成\textbf{“lktp”}或\textbf{反向钓鱼}操作，展现出极高的AI博弈能力。
    \item \textbf{指挥官面板}：前端使用 \textbf{React} 框架构建了\textbf{后台管理面板 (Admin Panel)}。作为操作员，可一键下达挖坑指令、动态调整Agent的\textbf{指令}、并实时监控“战果”。
\end{itemize}

\section{\makebox[\widthof{\faCogs}][c]{\color{CVBlue}\faCogs}\ 技术栈}
\begin{itemize}[nosep]
    \item \textbf{编程语言：} \textbf{Go}, Python, C++
    \item \textbf{开发工具：} SSH, Git, Vim, MakeFile,LaTex
    \item \textbf{操作系统：} Linux
\end{itemize}
\section{\makebox[\widthof{\faGraduationCap}][c]{\color{CVBlue}\faList}\ 获奖情况}
\begin{itemize}
	\item CC98年度用户 \hfill 2023.12
	
\end{itemize}
    
\section{\makebox[\widthof{\faInfo}][c]{\color{CVBlue}\faInfo}\ 其他}
\begin{itemize}[parsep=0.5ex]
    \item \textbf{技术博客：} \href{https://maksymilan.github.io/}{https://maksymilan.github.io/}
    \item \textbf{GitHub：} \href{https://github.com/maksymilan}{https://github.com/maksymilan} 
    \item \textbf{英语水平：} CET-4, CET-6
\end{itemize}
\end{document}